\documentclass[10pt,a4paper]{amsart}
\usepackage[margin=19mm]{geometry}
\usepackage{amsmath,amssymb,mathtools}
\usepackage[scr=rsfs,cal=euler]{mathalfa}
\usepackage{xfrac}
\usepackage[unicode,hidelinks,bookmarks=false]{hyperref}
\newcommand{\ie}{i.e.\ }
\newcommand{\cf}{cf.\ }
\newcommand{\resp}{respectively\ }
\newcommand{\Subsection}{\subsection}
\providecommand{\nobreakdash}{\nobreak\hbox{-}}
\setlength{\emergencystretch}{2em}
\multlinegap0pt
\usepackage{bm}
\usepackage{bbm}
\usepackage{dsfont}
\usepackage{xspace}
\usepackage{cancel}
\usepackage{mathtools}
\usepackage{amsthm}
\makeatletter
\@ifundefined{thm}{\newtheorem{thm}{Thm}}{}
\@ifundefined{claim}{\newtheorem{claim}{Claim}}{}
\@ifundefined{theorem}{\newtheorem{theorem}[thm]{Theorem}}{}
\makeatother
\title[Computability of Separation Axioms in Countable Second Count]{Computability of Separation Axioms in Countable Second Countable Spaces}
\author{Andrew DeLapo, David Gonzalez}
\date{Source: arXiv:2507.18564v1 (CC BY 4.0)}
\begin{document}
\maketitle

\noindent\emph{Source and licence.} Computability of Separation Axioms in Countable Second Countable Spaces. Version arXiv:2507.18564v1 (CC BY 4.0), distributed under the Creative Commons Attribution 4.0 International licence, as recorded in the arXiv licence metadata for this exact version and shown on the abstract page https://arxiv.org/abs/2507.18564v1. Full rights notice: licenses/arxiv-2507.18564-CC-BY-4.0.txt.\par\smallskip

\noindent\textit{Editorial scope.} Counted unit: the Theorem whose source label is \texttt{None} in the article (source0.tex lines 829--831, source label \texttt{None}), delivered together with the article's own complete original proof of it (source0.tex lines 833--965, one \texttt{proof} environment of 133 lines). No mathematics is written, repaired or re-derived: statement and proof are the article's own text line for line. Required context retained uncounted: thm:T3Completeness (lines 642--645). Every definition, display and label of those lines is retained unchanged. Omitted from this package and not redistributed: 1--641, 646--828, 832--832, 966--1481. Nothing inside a retained excerpt is omitted or altered: every retained body is delivered exactly as the article's lines are written, so no omission receipt of any kind accompanies this package. Citations: the retained excerpts carry no citation command at all, so no bibliography entry of the article is transcribed and no reference is redistributed. Reference adaptation: none was needed and none was made; every reference inside the delivered unit resolves inside it, so no unresolved reference or citation is produced. Printed slips kept literal: no printed word, symbol, hypothesis or number of the counted statement or of its proof was altered. Layout and class: the article's own class is \texttt{amsart} with packages including amssymb, xy, hyperref, color, pdfsync, graphicx, eucal, lmodern, cleveref, amsmath, amsthm, mathtools, MnSymbol, setspace; the wrapper is the lane's pinned \texttt{amsart} editorial wrapper at 10pt on A4, which restates the theorem-like environments the retained text needs and adds the article's own macro definitions that the retained text needs (none), with extra wrapper packages (bm, bbm, dsfont, xspace, cancel, mathtools). The delivered numbering is the unit's own. Assets: no retained span contains a figure environment, a graphics-inclusion command, a PDF-page inclusion, a table or a TikZ picture; no image file is copied into this package, the package declares no asset dependency and no image file is redistributed with it. Licence: version arXiv:2507.18564v1 of this article is distributed under the Creative Commons Attribution 4.0 International licence, as recorded in the arXiv licence metadata for this exact version and shown on the abstract page https://arxiv.org/abs/2507.18564v1. A full rights notice is delivered at licenses/arxiv-2507.18564-CC-BY-4.0.txt. Characters: every retained excerpt is pure ASCII, so the delivered engine, XeLaTeX with its default Latin Modern text font, reports no missing character.
\begin{theorem}\label{thm:T3Completeness}
    Let $T_3$-$CSC$ be the set of regular $CSC$ spaces.
    $T_3$-$CSC$ is $\Pi_5^0$ complete within within $T_2$-$CSC$.
\end{theorem}

    \begin{theorem}
        $T_{2\frac12}$-$CSC$ is $\Pi_5^0$ complete within within $T_2$-$CSC$.
    \end{theorem}

    \begin{proof}
        There are similarities with the proof techniques used in the proof of this theorem and those used for Theorem \ref{thm:T3Completeness}.
        That said, they are different enough that they require their own exposition.

         Our construction is built out of several basic modules of increasing complexity.
         Given $e$ coding a c.e.\ set $W_e$, we construct the following CSC space.
        We follow the convention that $\Phi_e$ enumerates at most one element into $W_e$ at each stage.
        Let $W_{e,s}^*$ be the computable set of stages at least $s$ where $W_e$ gets a new element.
        Let $W=\{2x,2x+1 \vert x\in W_{e,0}^*\}$.
        We construct the topology $\mathcal{V}_e=(V_i)_{i\in\omega}$ on $W\cup\{0\}$ where
    \[2x,2x+1\in V_{3s} \iff x\in W_{e,s}^*\cup\{0\}\]\[V_{3s+1}=\{2x\vert x\in W_{e,s+1}^*-W_{e,s}^*\}\]
    \[V_{3s+2}=\{2x+1\vert x\in W_{e,s+1}^*-W_{e,s}^*\}.\]
        It is useful notation to split $\mathcal{V}_e$ into three parts: $0$, $\mathcal{V}_e^0$ and $\mathcal{V}_e^1$.
        $0$ is just that; the element $0$.
        $\mathcal{V}_e^0$ consists of all of the non-zero, even elements of $\mathcal{V}_e$.
        $\mathcal{V}_e^1$ consists of all of the odd elements of $\mathcal{V}_e$.

    
    \begin{claim}\label{claim:basicUnit1/2}
        The element represented by $0$ in $\mathcal{V}_e$ is in $\overline{\mathcal{V}_e^0}$ and $\overline{\mathcal{V}_e^1}$ if and only if $W_e$ is finite.
    \end{claim}

    \begin{proof}
    If $W_e$ is finite, then there is some stage $s$ where $W_{e,s}^*$ is empty.
    In particular, $W_{e,s}^*\cup\{0\}=\{0\}$ isolates the point $0$.
    This means that $0$ is not in the closure of any open set that does not contain $0$, so in particular it is not in $\overline{\mathcal{V}_e^0}$ or $\overline{\mathcal{V}_e^1}$.

    If $W_e$ is infinite, then every open set around $0$ contains a subset of the form $V_{3s}$.
    Furthermore, each of the $V_{3s}$ contains infinitely many elements from both $\mathcal{V}_e^0$ and $\mathcal{V}_e^1$.
    If one were finite, observe that $|W_e|=|W_{e,s}\cup \mathcal{V}_e^i|$, and the right hand side would be finite, a contradiction to the assumption that $W_e$ is infinite.
    This means that every open set around $0$ contains elements from both $\mathcal{V}_e^0$ and $\mathcal{V}_e^1$.
    Or, what is the same, $0$ is in $\overline{\mathcal{V}_e^0}$ and $\overline{\mathcal{V}_e^1}$.
    \end{proof}


        It is worth explicitly noting that, by construction, all of the points that are not $0$ are always isolated in $\mathcal{V}_e$.

    Geometrically, $\mathcal{V}_e$ can be thought of as the point $0$ being placed at the center of $[-1,1]$ while the $n^{th}$ point enumerated into $W_e$ adds in the elements $\frac{1}{2^n}$ and $-\frac{1}{2^n}$.
    While $0$ is not a limit point from either direction after finitely many of these actions, it becomes a limit point from both directions after infinitely many.

    We now construct a larger basic unit of the construction.
    This unit will be made up of countably many $\mathcal{V}_e$ topologies and two additional points that we call $x^*$ and $y^*$.
    There will be closed neighborhoods separating $x^*$ and $y^*$, if and only if a $\Sigma_4^0$ condition is met.

    Let $A$ be a $\Sigma_4$ set and $f(a,c)$ be a function such that
    \[a\in A\iff \exists b\forall c\geq b ~ W_{f(a,c)} \text{ is finite.}\]
    Let $\mathcal{U}_{a,A}$ be the topology consisting of the basis for the topologies $\mathcal{V}_{f(a,c)}$ for each $c\in\omega$ and the the following additional open sets
    \[U_{2i}=\{x^*\}\cup \bigcup_{j\geq i}\mathcal{V}^0_{f(a,j)},\]
    \[U_{2i+1}=\{y^*\}\cup \bigcup_{j\geq i}\mathcal{V}^1_{f(a,j)}.\]
    In short, $\mathcal{U}_{a,A}$ is the disjoint union of the $\mathcal{V}_{f(a,c)}$ along with two extra points that are ``at the limit'' of the $\mathcal{V}^0_{f(a,c)}$ and $\mathcal{V}^1_{f(a,c)}$ respectively.

    \begin{claim}\label{claim:biggerUnit1/2}
        There are closed neighborhoods separating $x^*$ and $y^*$ if and only if $a\in A$.
    \end{claim}

    \begin{proof}
    Say that $a\in A$.
    Let $b'$ be the witness showing that $\forall c\geq b'~W_{f(a,c)} \text{ is finite.}$
    Consider $x^*\in U_{2b'}$ and $y^*\in U_{2b'+1}$.
    We claim that $\overline{U_{2b'}}=U_{2b'}$ and $\overline{U_{2b'+1}}=U_{2b'+1}$.
    We show only $\overline{U_{2b'}}=U_{2b'}$ as the other case is symmetric.
    Let $x\not\in U_{2b'}$, we must show that $x\not\in\overline{U_{2b'}}$, or what is the same, find an open set $V_x$ with $x\in V_x$ and $V_x\cap U_{2b'}=\emptyset$.
    Say that $x\in \mathcal{V}_{f(a,c)}$ with $c<b'$.
    The set $\mathcal{V}_{f(a,c)}$ is then itself such a $V_x$.
    On the other hand, say that $x\in \mathcal{V}_{f(a,c)}$ with $c\geq b'$.
    Because $x\not\in U_{2b'}$, this means that $x=0_c$ or $x\in \mathcal{V}_{f(a,c)}^1$ with $c\geq b'$.
    Because $W_{f(a,c)}$ is finite, this gives that, by Claim \ref{claim:basicUnit1/2}, $0_c$ is not in $\overline{U_{2b'}}$ and is isolated.
    If $x\in \mathcal{V}_{f(a,c)}^1$, the open set $U_{2b'+1}$ separates $x$ from $U_{2b'}$ as desired.
    As this covers all cases, this gives that $\overline{U_{2b'}}=U_{2b'}$ and symmetrically $\overline{U_{2b'+1}}=U_{2b'+1}$.
    As $x^*\in U_{2b'}$ yet $x^*\not\in U_{2b'+1}$ while $y^*\in U_{2b'+1}$ yet $y^*\not\in U_{2b'}$, this gives the desired separating closed neighborhoods.

    Say that $a\not\in A$.
    Consider open sets $X$ containing $x^*$ and $Y$ containing $y^*$.
    Any such $X$ must contain some $U_{2i}$ and $Y$ must contain some $U_{2j+1}$.
    We argue that $\overline{U_{2i}}\cap \overline{U_{2j+1}}\neq \emptyset$, which is enough to demonstrate the claim.
    Because $a\not\in A$ there is a $c\geq i,j$ such that $W_{f(a,c)}$ is infinite, and therefore $0_c$ is a limit point of both $\mathcal{V}_{f(a,c)}^0$ and $\mathcal{V}_{f(a,c)}^1$ by Claim \ref{claim:basicUnit1/2}.
    This means that $0_c\in \overline{\mathcal{V}_{f(a,c)}^0}\cap\overline{\mathcal{V}_{f(a,c)}^0}\subseteq\overline{U_{2i}}\cap \overline{U_{2j+1}},$ as desired.
    \end{proof}

    

    We now complete the construction of the desired space.
    Given a $\Pi_5^0$ set $B$ write
    \[n\in B \iff \forall m ~ n\in A_m,\]
    where $A_m$ is a $\Sigma_4$ condition.
    Let $\mathcal{Y}_n$ be the disjoint union of the spaces $\mathcal{U}_{A_m,n}$ over all values of $m$.
    We use $x^*_m$ and $y^*_m$ to denote the elements referred to in Claim \ref{claim:biggerUnit1/2} for the subspace $\mathcal{U}_{A_m,n}$.
    We use $0_{c,m}$ to denote the $c^{th}$ potential limit point in $\mathcal{U}_{A_m,n}$.
    In general, a subscript $m$ is added to the above-established notation to distinguish the different copies of $\mathcal{U}_{A_m,n}$.
    The reduction indicated in the theorem is given by the following claims.

    \begin{claim}\label{claim:T2for1/2}
        $\mathcal{Y}_n$ is $T_2$.
    \end{claim}

    \begin{proof}
        Consider two elements $x,y\in\mathcal{Y}_n$.
        The only (potentially) non-isolated points are the $x^*_m$, $y^*_m$, and the $0_{c,m}$.
        Any two isolated points are separated by the open sets only containing themselves.

        The opens $\mathcal{U}_{A_m,n}$ and $\mathcal{U}_{A_{m'},n}$ are disjoint when $m\neq m'$, so separate any elements coming from different pieces of the disjoint union.
        $x^*_m$ is separated from $0_{c,m}$ via the open sets $U_{2c+2,m}$ containing $x^*_m$ and $\mathcal{V}_{f(a,c,m)}$ containing $0_{c,m}$.
        $y^*_m$ is separated from $0_{c,m}$ via the open sets $U_{2c+3,m}$ containing $y^*_m$ and $\mathcal{V}_{f(a,c,m)}$ containing $0_{c,m}$.
        Similarly, $x^*_m$ and $y^*_m$ are separated from any other elements in the $\mathcal{V}_{f(a,c,m)}$.
        $x^*_m$ and $y^*_m$ are separated by the disjoint open sets $U_{0,m}$ and $U_{1,m}$.

        $0_{c,m}$ is similarly separated from other points not in the same copy of $\mathcal{V}_{f(a,c,m)}$.
        Lastly, $0_{c,m}$ is separated from other (isolated) points in $\mathcal{V}_{f(a,c,m)}$ by considering $V_{3s,c,m}$ where $s$ is larger than the index of the isolated point being considered.

        This covers all of the possible cases for pairs of points in $\mathcal{Y}_n$, showing that it is always Hausdorff. 
    \end{proof}

    \begin{claim}\label{claim:T21/2Reduction}
        $\mathcal{Y}_n$ is $T_{2\frac12}$ if and only if $n\in B$.
    \end{claim}

    \begin{proof}
        If $n\not\in B$ then for some $m$, $n\not\in A_m$.
        In particular, $x^*_m$ and $y^*_m$ fail to be separated by closed neighborhoods by Claim \ref{claim:biggerUnit1/2}.
        This means that $\mathcal{U}_{A_m,n}$, and indeed $\mathcal{Y}_n$, cannot be $T_{2\frac12}$.

        If $n\in B$, then for each $m$, $n\in A_m$.
        By Claim \ref{claim:biggerUnit1/2}, $x^*_m$ and $y^*_m$ are separated by closed neighborhoods for every $m$.
        By Claim \ref{claim:T2for1/2}, every element in $\mathcal{Y}_n$ is closed as a subset of $\mathcal{Y}_n$.
        This means that every isolated point is clopen.
        If $x$ is isolated, and $X$,$Y$ separate $x$ from $y$, $x$,$\overline{Y}$ are closed neighborhoods that separate $x$ from $y$.
        This is because $x\not\in \overline{Y}$ as $\{x\}$ is an open set disjoint from $Y$.
        Note that $0_{c,m}$ and $0_{c',m}$ for $c\neq c'$ are separated by the clopens $\mathcal{V}_{c,m}$ and $\mathcal{V}_{c',m}$.
        Furthermore, $0_{c,m}$ is separated from $x^*$ and $y^*$ by the clopens $\mathcal{V}_{c,m}$ and $U_{2c+2}$ or $U_{2c+3}$ respectively.
        Lastly, elements in different $A_m$ are separated by the clopens containing the entirety of the $A_m$ they are a part of.
    \end{proof}

    \end{proof}

\end{document}
